% =====================================================================
% Note on: Pricing Index-Linked Hybrids under a Local Volatility Market Model
% Mario Pucci (Banca IMI), 2012 — ssrn:2056277
% =====================================================================
\documentclass[11pt,a4paper]{article}

\newcommand{\papertag}{pucci\_2012\_hybrid}

% =====================================================================
% Shared preamble for paper notes
% Include with:  % =====================================================================
% Shared preamble for paper notes
% Include with:  % =====================================================================
% Shared preamble for paper notes
% Include with:  \input{../_common/preamble}
% =====================================================================

% --- Core packages ---
\usepackage[utf8]{inputenc}
\usepackage[T1]{fontenc}
\usepackage{lmodern}
\usepackage[margin=2.5cm]{geometry}
\usepackage{amsmath,amssymb,amsthm,mathtools}
\usepackage{bm}
\usepackage{booktabs}
\usepackage{array}
\usepackage{enumitem}
\usepackage{hyperref}
\usepackage{xcolor}
\usepackage{listings}
\usepackage{fancyhdr}
\usepackage{titlesec}

% --- Hyperref styling ---
\hypersetup{
  colorlinks=true,
  linkcolor=blue!50!black,
  urlcolor=blue!50!black,
  citecolor=blue!50!black
}

% --- Theorem-like environments ---
\theoremstyle{definition}
\newtheorem{defn}{Definition}
\newtheorem{remark}{Remark}
\newtheorem{example}{Example}

\theoremstyle{plain}
\newtheorem{prop}{Proposition}
\newtheorem{lemma}{Lemma}

% --- Coloured boxes for the structured sections ---
% Validation block, commentary, TODOs use coloured frames so the eye
% can find them when flipping through the note.
\usepackage[most]{tcolorbox}

\newtcolorbox{commentary}{
  colback=yellow!5,
  colframe=yellow!50!black,
  title=Commentary,
  fonttitle=\bfseries,
  breakable
}

\newtcolorbox{validation}{
  colback=green!3,
  colframe=green!50!black,
  title=Validation,
  fonttitle=\bfseries,
  breakable
}

\newtcolorbox{todobox}{
  colback=red!3,
  colframe=red!50!black,
  title=TODO / Open questions,
  fonttitle=\bfseries,
  breakable
}

% --- Code listing style for Python snippets in validation blocks ---
\lstdefinestyle{py}{
  language=Python,
  basicstyle=\ttfamily\footnotesize,
  keywordstyle=\color{blue!60!black},
  commentstyle=\color{gray},
  stringstyle=\color{red!60!black},
  showstringspaces=false,
  breaklines=true,
  frame=single,
  framesep=4pt,
  rulecolor=\color{gray!30}
}

% --- Page header: shows paper tag so a printed note is identifiable ---
\pagestyle{fancy}
\fancyhf{}
\fancyhead[L]{\small\textit{\papertag}}
\fancyhead[R]{\small\thepage}
\renewcommand{\headrulewidth}{0.4pt}

% --- Section spacing: tighter than article default ---
\titlespacing*{\section}{0pt}{1.5ex plus 0.5ex minus 0.2ex}{1ex plus 0.2ex}
\titlespacing*{\subsection}{0pt}{1.2ex plus 0.3ex minus 0.2ex}{0.6ex plus 0.2ex}

% =====================================================================

% --- Core packages ---
\usepackage[utf8]{inputenc}
\usepackage[T1]{fontenc}
\usepackage{lmodern}
\usepackage[margin=2.5cm]{geometry}
\usepackage{amsmath,amssymb,amsthm,mathtools}
\usepackage{bm}
\usepackage{booktabs}
\usepackage{array}
\usepackage{enumitem}
\usepackage{hyperref}
\usepackage{xcolor}
\usepackage{listings}
\usepackage{fancyhdr}
\usepackage{titlesec}

% --- Hyperref styling ---
\hypersetup{
  colorlinks=true,
  linkcolor=blue!50!black,
  urlcolor=blue!50!black,
  citecolor=blue!50!black
}

% --- Theorem-like environments ---
\theoremstyle{definition}
\newtheorem{defn}{Definition}
\newtheorem{remark}{Remark}
\newtheorem{example}{Example}

\theoremstyle{plain}
\newtheorem{prop}{Proposition}
\newtheorem{lemma}{Lemma}

% --- Coloured boxes for the structured sections ---
% Validation block, commentary, TODOs use coloured frames so the eye
% can find them when flipping through the note.
\usepackage[most]{tcolorbox}

\newtcolorbox{commentary}{
  colback=yellow!5,
  colframe=yellow!50!black,
  title=Commentary,
  fonttitle=\bfseries,
  breakable
}

\newtcolorbox{validation}{
  colback=green!3,
  colframe=green!50!black,
  title=Validation,
  fonttitle=\bfseries,
  breakable
}

\newtcolorbox{todobox}{
  colback=red!3,
  colframe=red!50!black,
  title=TODO / Open questions,
  fonttitle=\bfseries,
  breakable
}

% --- Code listing style for Python snippets in validation blocks ---
\lstdefinestyle{py}{
  language=Python,
  basicstyle=\ttfamily\footnotesize,
  keywordstyle=\color{blue!60!black},
  commentstyle=\color{gray},
  stringstyle=\color{red!60!black},
  showstringspaces=false,
  breaklines=true,
  frame=single,
  framesep=4pt,
  rulecolor=\color{gray!30}
}

% --- Page header: shows paper tag so a printed note is identifiable ---
\pagestyle{fancy}
\fancyhf{}
\fancyhead[L]{\small\textit{\papertag}}
\fancyhead[R]{\small\thepage}
\renewcommand{\headrulewidth}{0.4pt}

% --- Section spacing: tighter than article default ---
\titlespacing*{\section}{0pt}{1.5ex plus 0.5ex minus 0.2ex}{1ex plus 0.2ex}
\titlespacing*{\subsection}{0pt}{1.2ex plus 0.3ex minus 0.2ex}{0.6ex plus 0.2ex}

% =====================================================================

% --- Core packages ---
\usepackage[utf8]{inputenc}
\usepackage[T1]{fontenc}
\usepackage{lmodern}
\usepackage[margin=2.5cm]{geometry}
\usepackage{amsmath,amssymb,amsthm,mathtools}
\usepackage{bm}
\usepackage{booktabs}
\usepackage{array}
\usepackage{enumitem}
\usepackage{hyperref}
\usepackage{xcolor}
\usepackage{listings}
\usepackage{fancyhdr}
\usepackage{titlesec}

% --- Hyperref styling ---
\hypersetup{
  colorlinks=true,
  linkcolor=blue!50!black,
  urlcolor=blue!50!black,
  citecolor=blue!50!black
}

% --- Theorem-like environments ---
\theoremstyle{definition}
\newtheorem{defn}{Definition}
\newtheorem{remark}{Remark}
\newtheorem{example}{Example}

\theoremstyle{plain}
\newtheorem{prop}{Proposition}
\newtheorem{lemma}{Lemma}

% --- Coloured boxes for the structured sections ---
% Validation block, commentary, TODOs use coloured frames so the eye
% can find them when flipping through the note.
\usepackage[most]{tcolorbox}

\newtcolorbox{commentary}{
  colback=yellow!5,
  colframe=yellow!50!black,
  title=Commentary,
  fonttitle=\bfseries,
  breakable
}

\newtcolorbox{validation}{
  colback=green!3,
  colframe=green!50!black,
  title=Validation,
  fonttitle=\bfseries,
  breakable
}

\newtcolorbox{todobox}{
  colback=red!3,
  colframe=red!50!black,
  title=TODO / Open questions,
  fonttitle=\bfseries,
  breakable
}

% --- Code listing style for Python snippets in validation blocks ---
\lstdefinestyle{py}{
  language=Python,
  basicstyle=\ttfamily\footnotesize,
  keywordstyle=\color{blue!60!black},
  commentstyle=\color{gray},
  stringstyle=\color{red!60!black},
  showstringspaces=false,
  breaklines=true,
  frame=single,
  framesep=4pt,
  rulecolor=\color{gray!30}
}

% --- Page header: shows paper tag so a printed note is identifiable ---
\pagestyle{fancy}
\fancyhf{}
\fancyhead[L]{\small\textit{\papertag}}
\fancyhead[R]{\small\thepage}
\renewcommand{\headrulewidth}{0.4pt}

% --- Section spacing: tighter than article default ---
\titlespacing*{\section}{0pt}{1.5ex plus 0.5ex minus 0.2ex}{1ex plus 0.2ex}
\titlespacing*{\subsection}{0pt}{1.2ex plus 0.3ex minus 0.2ex}{0.6ex plus 0.2ex}

% =====================================================================
% House notation: shared symbols used across all paper notes.
% Each note imports this and may shadow individual macros if a paper
% uses a clashing convention — but the *default* meaning is fixed here.
%
% Include with:  % =====================================================================
% House notation: shared symbols used across all paper notes.
% Each note imports this and may shadow individual macros if a paper
% uses a clashing convention — but the *default* meaning is fixed here.
%
% Include with:  % =====================================================================
% House notation: shared symbols used across all paper notes.
% Each note imports this and may shadow individual macros if a paper
% uses a clashing convention — but the *default* meaning is fixed here.
%
% Include with:  \input{../_common/notation}
% =====================================================================

% --- Measures and expectations ---
\newcommand{\Q}{\mathbb{Q}}              % risk-neutral measure
\newcommand{\Qt}{\mathbb{Q}^{T}}         % T-forward measure
\newcommand{\E}{\mathbb{E}}              % expectation
\newcommand{\Et}[1]{\E^{#1}}             % expectation under measure #1
\newcommand{\Ftime}[1]{\mathcal{F}_{#1}} % filtration at #1

% --- Discount and forward objects ---
\newcommand{\Disc}[2]{D_{#1,#2}}         % discount factor D_{t,T}
\newcommand{\Bond}[2]{P_{#1}(#2)}        % bond price P_t(y)
\newcommand{\Ann}[1]{A_{#1}}             % annuity / level
\newcommand{\Numer}[1]{N_{#1}}           % numeraire
\newcommand{\Fwd}[2]{F_{#1,#2}}          % forward price F_{t,T}
\newcommand{\Fwdpx}[2]{\mathrm{Fwd}_{#1}(#2)} % verbose forward-price F_t(T)

% --- Rates ---
\newcommand{\IRR}{\mathrm{IRR}}          % internal rate of return
\newcommand{\Yld}{y}                     % bond yield variable
\newcommand{\Strike}{K}                  % strike (price-space)
\newcommand{\Lock}{L}                    % locked yield / rate
\newcommand{\Repo}{r^{\mathrm{repo}}}    % repo rate
\newcommand{\Fund}{r^{\mathrm{fun}}}     % unsecured funding rate
\newcommand{\OIS}{r^{\mathrm{ois}}}      % OIS rate
\newcommand{\Coll}{r^{\mathrm{c}}}       % collateral rate
\newcommand{\Rcmt}{R^{\mathrm{cmt}}}     % constant-maturity-treasury rate
\newcommand{\Rswp}{R^{\mathrm{swp}}}     % swap rate (used for risky variant)
\newcommand{\Rec}{\mathrm{Rec}}          % recovery rate
\newcommand{\NPV}{\mathrm{NPV}}          % present-value functional
\newcommand{\PVone}{\mathrm{PV01}}       % risky annuity
\newcommand{\Loss}{\mathrm{Loss}}        % default-leg integral
\newcommand{\Sasw}{S^{\mathrm{ASW}}}     % asset-swap spread
\newcommand{\Scds}{S^{\mathrm{CDS}}}     % CDS spread
\newcommand{\Basis}{\mathrm{Basis}}      % CDS-bond basis

% --- Sensitivities ---
\newcommand{\Diff}[2]{D_{#1}^{#2}}       % Euler's notation: D_x^k[f]
\newcommand{\Riskf}{\mathrm{RiskFactor}} % bond risk factor (= -DV01*1e4)
\newcommand{\DVone}{\mathrm{DV01}}

% --- Indicator / misc ---
\newcommand{\Ind}{\mathbf{1}}
\newcommand{\given}{\,|\,}
\newcommand{\dd}{\,\mathrm{d}}

% --- Greeks-style ---
\newcommand{\Gam}{\Gamma}
\newcommand{\Vega}{\mathcal{V}}

% =====================================================================

% --- Measures and expectations ---
\newcommand{\Q}{\mathbb{Q}}              % risk-neutral measure
\newcommand{\Qt}{\mathbb{Q}^{T}}         % T-forward measure
\newcommand{\E}{\mathbb{E}}              % expectation
\newcommand{\Et}[1]{\E^{#1}}             % expectation under measure #1
\newcommand{\Ftime}[1]{\mathcal{F}_{#1}} % filtration at #1

% --- Discount and forward objects ---
\newcommand{\Disc}[2]{D_{#1,#2}}         % discount factor D_{t,T}
\newcommand{\Bond}[2]{P_{#1}(#2)}        % bond price P_t(y)
\newcommand{\Ann}[1]{A_{#1}}             % annuity / level
\newcommand{\Numer}[1]{N_{#1}}           % numeraire
\newcommand{\Fwd}[2]{F_{#1,#2}}          % forward price F_{t,T}
\newcommand{\Fwdpx}[2]{\mathrm{Fwd}_{#1}(#2)} % verbose forward-price F_t(T)

% --- Rates ---
\newcommand{\IRR}{\mathrm{IRR}}          % internal rate of return
\newcommand{\Yld}{y}                     % bond yield variable
\newcommand{\Strike}{K}                  % strike (price-space)
\newcommand{\Lock}{L}                    % locked yield / rate
\newcommand{\Repo}{r^{\mathrm{repo}}}    % repo rate
\newcommand{\Fund}{r^{\mathrm{fun}}}     % unsecured funding rate
\newcommand{\OIS}{r^{\mathrm{ois}}}      % OIS rate
\newcommand{\Coll}{r^{\mathrm{c}}}       % collateral rate
\newcommand{\Rcmt}{R^{\mathrm{cmt}}}     % constant-maturity-treasury rate
\newcommand{\Rswp}{R^{\mathrm{swp}}}     % swap rate (used for risky variant)
\newcommand{\Rec}{\mathrm{Rec}}          % recovery rate
\newcommand{\NPV}{\mathrm{NPV}}          % present-value functional
\newcommand{\PVone}{\mathrm{PV01}}       % risky annuity
\newcommand{\Loss}{\mathrm{Loss}}        % default-leg integral
\newcommand{\Sasw}{S^{\mathrm{ASW}}}     % asset-swap spread
\newcommand{\Scds}{S^{\mathrm{CDS}}}     % CDS spread
\newcommand{\Basis}{\mathrm{Basis}}      % CDS-bond basis

% --- Sensitivities ---
\newcommand{\Diff}[2]{D_{#1}^{#2}}       % Euler's notation: D_x^k[f]
\newcommand{\Riskf}{\mathrm{RiskFactor}} % bond risk factor (= -DV01*1e4)
\newcommand{\DVone}{\mathrm{DV01}}

% --- Indicator / misc ---
\newcommand{\Ind}{\mathbf{1}}
\newcommand{\given}{\,|\,}
\newcommand{\dd}{\,\mathrm{d}}

% --- Greeks-style ---
\newcommand{\Gam}{\Gamma}
\newcommand{\Vega}{\mathcal{V}}

% =====================================================================

% --- Measures and expectations ---
\newcommand{\Q}{\mathbb{Q}}              % risk-neutral measure
\newcommand{\Qt}{\mathbb{Q}^{T}}         % T-forward measure
\newcommand{\E}{\mathbb{E}}              % expectation
\newcommand{\Et}[1]{\E^{#1}}             % expectation under measure #1
\newcommand{\Ftime}[1]{\mathcal{F}_{#1}} % filtration at #1

% --- Discount and forward objects ---
\newcommand{\Disc}[2]{D_{#1,#2}}         % discount factor D_{t,T}
\newcommand{\Bond}[2]{P_{#1}(#2)}        % bond price P_t(y)
\newcommand{\Ann}[1]{A_{#1}}             % annuity / level
\newcommand{\Numer}[1]{N_{#1}}           % numeraire
\newcommand{\Fwd}[2]{F_{#1,#2}}          % forward price F_{t,T}
\newcommand{\Fwdpx}[2]{\mathrm{Fwd}_{#1}(#2)} % verbose forward-price F_t(T)

% --- Rates ---
\newcommand{\IRR}{\mathrm{IRR}}          % internal rate of return
\newcommand{\Yld}{y}                     % bond yield variable
\newcommand{\Strike}{K}                  % strike (price-space)
\newcommand{\Lock}{L}                    % locked yield / rate
\newcommand{\Repo}{r^{\mathrm{repo}}}    % repo rate
\newcommand{\Fund}{r^{\mathrm{fun}}}     % unsecured funding rate
\newcommand{\OIS}{r^{\mathrm{ois}}}      % OIS rate
\newcommand{\Coll}{r^{\mathrm{c}}}       % collateral rate
\newcommand{\Rcmt}{R^{\mathrm{cmt}}}     % constant-maturity-treasury rate
\newcommand{\Rswp}{R^{\mathrm{swp}}}     % swap rate (used for risky variant)
\newcommand{\Rec}{\mathrm{Rec}}          % recovery rate
\newcommand{\NPV}{\mathrm{NPV}}          % present-value functional
\newcommand{\PVone}{\mathrm{PV01}}       % risky annuity
\newcommand{\Loss}{\mathrm{Loss}}        % default-leg integral
\newcommand{\Sasw}{S^{\mathrm{ASW}}}     % asset-swap spread
\newcommand{\Scds}{S^{\mathrm{CDS}}}     % CDS spread
\newcommand{\Basis}{\mathrm{Basis}}      % CDS-bond basis

% --- Sensitivities ---
\newcommand{\Diff}[2]{D_{#1}^{#2}}       % Euler's notation: D_x^k[f]
\newcommand{\Riskf}{\mathrm{RiskFactor}} % bond risk factor (= -DV01*1e4)
\newcommand{\DVone}{\mathrm{DV01}}

% --- Indicator / misc ---
\newcommand{\Ind}{\mathbf{1}}
\newcommand{\given}{\,|\,}
\newcommand{\dd}{\,\mathrm{d}}

% --- Greeks-style ---
\newcommand{\Gam}{\Gamma}
\newcommand{\Vega}{\mathcal{V}}


\title{Note on: \emph{Pricing Index-Linked Hybrids under a Local Volatility Market Model}\\
       \large Mario Pucci (Banca IMI), 2012}
\author{Reading notes}
\date{\today}

\begin{document}
\maketitle

% =====================================================================
\section*{Metadata}
\begin{tabular}{@{}p{3.5cm}p{11cm}@{}}
\toprule
Title           & Pricing Index-Linked Hybrids under a Local Volatility Market Model \\
Author          & Mario Pucci (Banca IMI, Milan) \\
Date            & 22 February 2012 (revised May 2012) \\
Source          & \href{https://ssrn.com/abstract=2056277}{ssrn:2056277} \\
Local file      & \texttt{papers/pucci\_2012\_hybrid/ssrn2056277.pdf} \\
Tags            & \texttt{hybrid}, \texttt{equity-rates}, \texttt{cash-swaption},
                  \texttt{local-vol}, \texttt{cms}, \texttt{forward-measure} \\
Date read       & \today \\
\bottomrule
\end{tabular}

% =====================================================================
\section*{One-paragraph summary}
A short note proposing a ``market model'' for pricing a cash-settled
swaption whose strike depends on the level of an equity index at the
exercise date $T$: payoff
$\widehat A_T \cdot (\theta(S_T - U_T))^+$ where $\widehat A_T$ is the
cash annuity, $S_T$ the forward swap rate, and $U_T$ the index
forward. The key trick is to switch from the swaption measure to the
$T$-forward measure: the cash annuity then becomes a deterministic
function of the convexity-adjusted forward swap rate $F_{T T} = S_T$,
so only the joint distribution of $(F, U)$ matters. Both $F$ and $U$
are modelled as local-volatility martingales under the $T$-forward
measure (calibratable individually to vanilla swaption and vanilla
index option markets), with an exogenously specified correlation
$\rho$. The framework keeps every dynamics ``market observable'':
$F_{0T}$ from CMS quotes / static replication, $\sigma_F$ from
swaptions, $\sigma_U$ from index options, $\rho$ from historical
data.

% =====================================================================
\section{Setup and notation}

\begin{tabular}{@{}p{3cm}p{3.5cm}p{8.5cm}@{}}
\toprule
Paper                   & House                       & Meaning \\
\midrule
$U_t$                   & ---                         & Spot index level \\
$U_{tT}$                & ---                         & Forward price of the index at $t$ for delivery $T$ \\
$S_t$                   & $R^{swp}_t$                 & Forward swap rate at $t$ for tenor $T_1, \ldots, T_n$ \\
$A_t$                   & $\Ann{t}$                   & Risk-free annuity, $\sum_i \alpha_i \Disc{t}{T_i}$ \\
$\widehat A_T$          & $\widehat{\Ann{T}}$         & Cash annuity (flat-curve discount on $S_T$) \\
$\theta$                & $\theta \in \{+1, -1\}$     & Call ($+1$) / put ($-1$) flag \\
$T = T_0$               & $T$                         & Swaption exercise date \\
$T_1, \ldots, T_n$      & $T_i$                       & Underlying swap payment dates \\
$D_{tT_i}$              & $\Disc{t}{T_i}$             & Risk-free DF \\
$F_{tT}$                & $F_{tT}$                    & $\mathbb{E}^T_t[S_T]$, convexity-adjusted forward swap rate \\
$\sigma_F, \sigma_U$    & ---                         & Local volatility functions \\
$W^F, W^U$              & ---                         & Driving BMs under the $T$-forward measure \\
$\rho$                  & $\rho$                      & Instantaneous correlation \\
$K$                     & $K$                         & Vanilla swaption / index option strike \\
\bottomrule
\end{tabular}

% =====================================================================
\section{Key equations}

\paragraph{Cash annuity (flat-curve proxy):}
\begin{equation}\label{eq:hybrid-cash-ann}
  \widehat A_T \;=\; \sum_{i=1}^{n} \frac{\alpha_i}{\bigl(1 + \alpha_i S_T\bigr)^{\mathrm{yf}(T, T_i)}},
\end{equation}
with $\alpha_i = \mathrm{yf}(T_{i-1}, T_i)$.
\textit{Says:} the cash annuity replaces market discount factors with
synthetic ones computed from a flat curve at $S_T$. Standard for
cash-settled swaptions in the European market. By construction,
$\widehat A_T$ is a function of $S_T$ alone.

\paragraph{Index-linked cash-settled swaption payoff:}
\begin{equation}\label{eq:hybrid-payoff}
  \widehat A_T \, \bigl(\theta(S_T - U_T)\bigr)^+.
\end{equation}
\textit{Says:} a cash-settled swaption with strike replaced by the
realised level of the index $U_T$.

\paragraph{Generic numeraire-pair PV:}
\begin{equation}\label{eq:hybrid-pv-generic}
  V_t \;=\; \Numer{t} \, \mathbb{E}^N_t\!\left[ \theta \, \widehat A_T \, (\theta(S_T - U_T))^+ / \Numer{T} \right].
\end{equation}

\paragraph{Switch to $T$-forward measure:}
With $F_{tT} \equiv \mathbb{E}^T_t[S_T]$ the convexity-adjusted
forward swap rate and noting $F_{TT} = S_T$,
\begin{equation}\label{eq:hybrid-pv-fwd}
  V_t \;=\; \Disc{t}{T}\, \mathbb{E}^T_t\!\left[ \widehat A_T \, \bigl(\theta(F_{TT} - U_{TT})\bigr)^+ \right].
\end{equation}
\textit{Says:} the change of numeraire reduces the joint problem to
specifying the bivariate distribution of $(F_{TT}, U_{TT})$ under the
$T$-forward measure; no need to model $\widehat A_T / A_T$ separately
because $\widehat A_T$ is a function of $S_T = F_{TT}$.

\paragraph{Vanilla calibration targets} (priced under the same model):

vanilla swaption (cash-settled):
\begin{equation}\label{eq:hybrid-vanilla-swp}
  V^{\text{swp}}_t \;=\; \Disc{t}{T}\, \mathbb{E}^T_t\!\left[ \widehat A_T \, \bigl(\theta(F_{TT} - K)\bigr)^+ \right],
\end{equation}

CMS cap / floor on the convexity-adjusted rate (often not liquid):
\begin{equation}\label{eq:hybrid-vanilla-cms}
  V^{\text{cms}}_t \;=\; \Disc{t}{T}\, \mathbb{E}^T_t\!\left[ \bigl(\theta(F_{TT} - K)\bigr)^+ \right],
\end{equation}

vanilla index option:
\begin{equation}\label{eq:hybrid-vanilla-idx}
  V^{\text{idx}}_t \;=\; \Disc{t}{T}\, \mathbb{E}^T_t\!\left[ \bigl(\theta(U_{TT} - K)\bigr)^+ \right].
\end{equation}

\paragraph{Local-volatility dynamics under $T$-forward measure:}
\begin{align}
  \frac{\dd F_{tT}}{F_{tT}} &= \sigma_F(t, F_{tT})\, \dd (W^F)^T_t, \label{eq:hybrid-dyn-F}\\
  \frac{\dd U_{tT}}{U_{tT}} &= \sigma_U(t, U_{tT})\, \dd (W^U)^T_t, \label{eq:hybrid-dyn-U}\\
  \dd\langle (W^F)^T, (W^U)^T\rangle_t &= \rho\, \dd t.
\end{align}
\textit{Says:} both forwards are martingales under $\mathbb{Q}^T$,
each with its own Dupire-style local-volatility function. The
correlation $\rho$ is exogenous to the calibration of marginals.

% =====================================================================
\section{Derivations \& commentary}

\begin{commentary}
The structural simplification is that $\widehat A_T$ is, by
definition \eqref{eq:hybrid-cash-ann}, a deterministic function of
$S_T$. So switching to $T$-forward leaves a problem entirely in
$(F_{TT}, U_{TT})$. Compared with traditional hybrid approaches
(short-rate + equity, e.g.\ G2++/BS or Heston-HW), this is much
lighter: no unobservable short-rate factors, no joint
stochastic-volatility specification, no FX-style hybrid correlation
matrix.
\end{commentary}

\begin{commentary}
The trade-off is honesty about correlation. $\rho$ between
$F_{tT}$ and $U_{tT}$ is not directly observable: there is no
liquid product whose price depends only on it. The paper's
suggestions are (i) interpolate forward swap rates along zero curves
of an index series as a proxy for the convexity-adjusted rate, then
estimate historical $\rho$ with $U_{tT}$; or (ii) use Black's ATM
approximation on historical forward swap rates to back out a
historical $F_{tT}$ time series. Both are heuristics. As with the
CMASW paper, the model honestly admits that correlation is the soft
spot.
\end{commentary}

\begin{commentary}
The convexity-adjusted CMS rate $F_{0T}$ is treated as a market
input. Practically, this is read off CMS spread quotes or a CMS
interpolator (often SABR-based) on the rates desk. The note's
calibration recipe is therefore:
\begin{enumerate}[noitemsep]
  \item $F_{0T}$ from CMS market.
  \item $\sigma_F(\cdot, \cdot)$ from swaption smile via
        \eqref{eq:hybrid-vanilla-swp} (Dupire on the
        cash-swaption-implied marginal of $F$).
  \item $\sigma_U(\cdot, \cdot)$ from index option smile via
        \eqref{eq:hybrid-vanilla-idx} (standard Dupire on the index).
  \item $\rho$ exogenous, ideally historical.
\end{enumerate}
\end{commentary}

\begin{commentary}
The paper says the framework extends to ``other multi-asset hybrid
payoffs of European type''. The key ingredients are: (i) the payoff
must be European; (ii) the cash component must reduce to a function
of a single underlying after the measure change; (iii) the remaining
underlyings need vanilla-option markets for marginal calibration.
Cliquets, autocallables and other path-dependent hybrids are out of
scope.
\end{commentary}

% =====================================================================
\section{Validation}

\begin{validation}
This is a \emph{specification} of mathematical objects the paper
defines and the numerical values they should take. It says nothing
about how those objects are implemented or invoked.

\textbf{Quantities the paper defines:}
\begin{itemize}[noitemsep]
  \item Cash annuity $\widehat A_T$ via \eqref{eq:hybrid-cash-ann},
        a deterministic function of $S_T$.
  \item Index-linked cash-swaption payoff at $T$ via
        \eqref{eq:hybrid-payoff}.
  \item $T$-forward measure pricing formula
        \eqref{eq:hybrid-pv-fwd}; convexity-adjusted forward
        $F_{tT} = \mathbb{E}^T_t[S_T]$.
  \item Three vanilla-option pricing formulas in the same framework
        \eqref{eq:hybrid-vanilla-swp}--\eqref{eq:hybrid-vanilla-idx},
        used as calibration targets.
  \item Joint local-volatility dynamics
        \eqref{eq:hybrid-dyn-F}--\eqref{eq:hybrid-dyn-U} with
        instantaneous correlation $\rho$.
\end{itemize}

\textbf{Canonical numerical case}: the paper provides no worked
numerical example. Suggested test bench: a 1Y expiry on a 10Y
cash-settled swaption (10Y tenor, annual payments) struck against a
major equity index (e.g.\ S\&P 500 or EuroStoxx 50). Use $F_{0T}$
read off a CMS interpolator, swaption ATM vol $\sim 30\%$, index ATM
vol $\sim 20\%$, $\rho \in \{-0.3, 0.0, +0.3\}$ as a sensitivity grid.

\textbf{Expected numerical values}: none printed in the paper.

\textbf{Equation $\leftrightarrow$ quantity map:}

\smallskip
\begin{tabular}{@{}p{2.8cm}p{3.8cm}p{6.2cm}@{}}
\toprule
Paper eq.                          & Symbol                  & Quantity \\
\midrule
\eqref{eq:hybrid-cash-ann}         & $\widehat A_T$          & cash annuity as deterministic fn of $S_T$ \\
\eqref{eq:hybrid-payoff}           & $\widehat A_T (S_T - U_T)^+$ & index-linked cash-swaption payoff \\
\eqref{eq:hybrid-pv-fwd}           & $V_t$                   & price under $T$-forward measure \\
\eqref{eq:hybrid-vanilla-swp}      & $V^{\text{swp}}_t$      & vanilla cash-swaption price (calibration target) \\
\eqref{eq:hybrid-vanilla-cms}      & $V^{\text{cms}}_t$      & CMS cap/floor price (alt.\ calibration target) \\
\eqref{eq:hybrid-vanilla-idx}      & $V^{\text{idx}}_t$      & vanilla index option price (calibration target) \\
\eqref{eq:hybrid-dyn-F}--\eqref{eq:hybrid-dyn-U} & $F_{tT}, U_{tT}$ & local-vol dynamics under $T$-forward \\
\bottomrule
\end{tabular}
\smallskip

\textbf{Invariants and identities to verify:}
\begin{enumerate}[noitemsep]
  \item \emph{Cash annuity at par}: $\widehat A_T$ evaluated at
        $S_T = R^{swp}_0$ (i.e.\ ATM forward) should match the level
        consistent with the market discount curve to within a
        first-order correction; the gap is the cash-vs-market annuity
        difference standard in cash-settled-swaption pricing.
  \item \emph{Recovery of vanilla prices}: pricing
        \eqref{eq:hybrid-vanilla-swp} via the calibrated $\sigma_F$
        reproduces the input swaption smile by construction
        (Dupire); same for \eqref{eq:hybrid-vanilla-idx}.
  \item \emph{Forward consistency}: $F_{tT}$ is a $\mathbb{Q}^T$-martingale
        by construction, so $\mathbb{E}^T_t[F_{TT}] = F_{tT}$. Same
        for $U_{tT}$.
  \item \emph{Limit at $\rho = 0$}: with independent $F$ and $U$,
        the hybrid price reduces to an iterated expectation over
        their marginals, computable in 1D once $\sigma_F, \sigma_U$
        are pinned.
  \item \emph{Strike intermediate values}: replacing $U_T$ in the
        payoff by a constant $K$ reduces \eqref{eq:hybrid-pv-fwd} to
        \eqref{eq:hybrid-vanilla-swp}, the vanilla cash-swaption,
        and the hybrid model and vanilla model must agree.
  \item \emph{Convexity adjustment consistency}: at $t = 0$, the
        difference between $F_{0T}$ (the convexity-adjusted forward)
        and the spot forward swap rate $S_0$ should equal the
        Hagan / Pelsser convexity adjustment under matching
        assumptions.
\end{enumerate}
\end{validation}

% =====================================================================
\section{Glossary of symbols (paper's notation)}
\begin{tabular}{@{}p{3cm}p{12cm}@{}}
\toprule
$U_t$                   & Spot equity index level at $t$ \\
$U_{tT}$                & Forward price of $U$ at $t$ for delivery $T$ \\
$S_t$                   & Forward swap rate at $t$ with payment schedule $T_1, \ldots, T_n$ \\
$A_t$                   & Risk-free annuity \\
$\widehat A_T$          & Cash annuity (flat-curve discount on $S_T$) \\
$\theta$                & $+1$ payer / $-1$ receiver flag \\
$T = T_0$               & Swaption exercise date \\
$T_1, \ldots, T_n$      & Underlying swap payment dates \\
$D_{tT}$                & Risk-free discount factor \\
$F_{tT}$                & Convexity-adjusted forward swap rate, $\mathbb{E}^T_t[S_T]$ \\
$\mathbb{Q}^T$          & $T$-forward measure (numeraire $D_{tT}$) \\
$\sigma_F, \sigma_U$    & Local-volatility functions for $F$ and $U$ \\
$(W^F)^T, (W^U)^T$      & Driving Brownian motions under $\mathbb{Q}^T$ \\
$\rho$                  & Instantaneous correlation of $(W^F)^T$ and $(W^U)^T$ \\
$K$                     & Strike for vanilla calibration instruments \\
\bottomrule
\end{tabular}

% =====================================================================
\begin{todobox}
\begin{itemize}[noitemsep]
  \item Build a worked numerical example since the paper provides
        none: 1Y$\times$10Y cash-swaption with strike $= U_T$ where
        $U$ is a major equity index; benchmark against an iterated
        BS-style approximation as a sanity check.
  \item Calibration pipeline: implement steps (i)--(iv) on a synthetic
        market (flat smiles, synthetic CMS rates) to verify that
        vanilla prices are recovered exactly.
  \item Correlation sensitivity: produce a price-vs-$\rho$ curve at
        plausible ATM vols ($\sigma_F, \sigma_U \sim 20\%-30\%$) to
        gauge model risk on the unobservable input.
  \item Compare with a flat-vol Black approximation
        $V \approx D_{0T} \widehat A_0 \cdot \text{BS}(F_0, U_0, \tilde\sigma, T)$
        with $\tilde\sigma^2 = \sigma_F^2 - 2\rho\sigma_F\sigma_U + \sigma_U^2$.
\end{itemize}
\end{todobox}

\end{document}
